\section{Introduction}
\label{sec:intro}

Autism spectrum disorder (ASD) is a neurodevelopmental condition affecting
roughly 1--2\,\% of the global population, with twin studies placing
heritability near 80\,\%~\cite{havdahl2021}. ASD's genetic architecture is
markedly polygenic: hundreds of common variants of small effect act
alongside rarer high-impact mutations, with a particular concentration in
genes regulating prenatal brain development and synaptic
function~\cite{satterstrom2020,grove2019}.

Genome-wide association studies (GWAS) have flagged more than one hundred
loci associated with ASD, but each individual variant explains only a
vanishing fraction of risk. Polygenic risk scores aggregate these tiny
signals into a per-individual estimate of genetic liability and have
become the dominant framework for translating GWAS evidence into a
quantitative predictor. Despite ever-larger sample sizes, however, ASD
PRS still account for less than 5\,\% of risk
variance~\cite{grove2019,havdahl2021} -- a striking gap given the trait's
heritability.

Two construction choices contribute to this shortfall. First, the
canonical pruning-and-thresholding (P+T) procedure retains only variants
clearing a stringent significance bar, deliberately discarding the bulk of
the genome's distributed signal that arises from a highly polygenic trait.
Second, even more sophisticated Bayesian shrinkage methods such as
PRS-CS~\cite{ge2019prscs} and LDpred2~\cite{prive2020ldpred2} retain the
additive linear assumption and treat variants as exchangeable conditional
on local LD, ignoring the rich biological structure -- shared pathways,
co-expression, protein--protein interactions -- that demonstrably links
ASD risk genes~\cite{satterstrom2020,dong2023}.

These two limitations motivate the present work. Rather than incrementally
collecting more samples, we ask whether a method that explicitly
incorporates SNP-level functional annotations and gene--gene interaction
networks can extract additional predictive signal from the
\emph{existing} ASD summary-statistics releases. We do so on
summary-statistics-only inputs, sidestepping the long timelines and
restricted-access requirements of individual-level genotype cohorts.

\subsection*{Contributions}

\begin{itemize}[leftmargin=1.2em,nosep]
  \item A reproducible sumstats-only pipeline that processes the
        iPSYCH--PGC and SPARK + iPSYCH + PGC ASD releases end-to-end,
        producing per-SNP weight files and an evaluation table for every
        method.
  \item An algebraic derivation that yields an independent SPARK-only
        Z-score per SNP from the two public sumstats files, providing a
        clean held-out evaluation target without requiring access to
        SPARK individual genotypes.
  \item Two ML-enhanced PRS strategies -- (i) feature-based gradient
        boosting and an MLP re-weighter, and (ii) a graph-neural-network
        re-weighter operating on a gene--gene interaction graph -- both
        compared head-to-head with P+T and a PRS-CS-style Bayesian
        baseline on a chromosomal hold-out.
  \item An enrichment analysis quantifying the alignment of each method's
        top-decile gene scores with the SFARI Gene catalogue and
        relevant Gene Ontology (GO) and KEGG pathways.
\end{itemize}
